\section{ Introduction}


Extraction of  parton distribution functions (PDFs) $f(x)$ \cite{Feynman:1973xc}
on the lattice is a challenging problem attracting a lot of attention.
 The usual  method to approach PDFs on the lattice  is to calculate  their moments.
  However, recently, X. Ji  \cite{Ji:2013dva} suggested a method
 allowing a calculation  of  PDFs   as  functions of $x$.

 Since the PDFs
are  related to   matrix elements
of bilocal operators on   the light cone $z^2=0$, this was a stumbling block preventing
a  direct  calculation of these functions in  the   lattice
 gauge  theory
 formulated in  Euclidean space.

  To overcome this difficulty, X. Ji proposes to use purely space-like
separations $z=(0,0,0,z_3)$.
The functions in this case are   quasi-PDFs  $Q(y,p_3)$ describing the distribution
of   the $p_3$  hadron  momentum  component.
The key point is that quasi-PDFs $Q(y,p_3)$
   tend to  usual PDFs $f(y)$
     in the \mbox{$p_3 \to \infty$}  limit.
     The same method can be applied to distribution amplitudes (DAs).
The results of  quasi-PDF calculations on the  lattice
   were reported
in Refs.  \cite{Lin:2014zya,Chen:2016utp,Alexandrou:2015rja}
and of the pion quasi-DA   in Ref. \cite{Zhang:2017bzy}.



Recent papers \cite{Radyushkin:2016hsy,Radyushkin:2017gjd}
by one of the authors (A.R.) contain an investigation of the  nonperturbative $p_3$-evolution of
quasi-PDFs and quasi-DAs. This study  is based on
 the formalism of  virtuality distribution functions
 \cite{Radyushkin:2014vla,Radyushkin:2015gpa}.
 The approach developed in Refs.
 \cite{Radyushkin:2016hsy,Radyushkin:2017gjd}  has  established  a connection between
the   quasi-PDFs  and
  the   ``straight-link'' transverse momentum dependent
 distributions    (TMDs) ${\cal F} (x, k_\perp^2)$.
 Starting  from  simple models for TMDs,  models were
 built  for the    nonperturbative evolution of
quasi-PDFs.  The derived curves agree qualitatively with the patterns of
\mbox{$p_3$-evolution}  produced by lattice simulations.

The structure  of quasi-PDFs was further studied in Ref.   \cite{Radyushkin:2017cyf}.
 It was   shown  that,
when a  hadron is moving, the parton $k_3$ momentum
may be treated  as
coming  from two sources. The hadron's motion as a whole yields
the  $xp_3$  part,  which is governed by the dependence of
 the TMD  ${\cal F} (x, \kappa^2)$ on its first argument namely $x$.
  The residual part $k_3-xp_3$
is controlled by the  way that the TMD depends on its second argument,
$\kappa^2$,
which dictates the shape of the primordial
 rest-frame
 momentum distribution.    Quasi-PDFs due to their convolution nature possess a rather  involved pattern of their
 \mbox{$p_3$-evolution,} making mandatory relatively big values $p_3 \gtrsim 3$ GeV in order to safely approach the PDF limit.

To accelerate the convergence,   a different  approach for the
PDF extraction from lattice  calculations
 was  proposed  \cite{Radyushkin:2017cyf}. It is  based on the concept of
{\it pseudo-PDFs} $ {\cal P} (x, z_3^2) $.  They  generalize
the light-cone PDFs $f(x)$ onto spacelike intervals like
 $z=(0,0,0,z_3)$.   The pseudo-PDFs are Fourier transforms
 of the  \mbox{{\it Ioffe-time}  \cite{Ioffe:1969kf}}  \mbox{{\it distributions} \cite{Braun:1994jq}}
${\cal M} (\nu, z_3^2)$
 which are generically given by matrix elements  $  \langle p |   \phi(0) \phi (z)|p \rangle $
 written as  functions of $\nu = p_3 z_3$ and $z_3^2$.
In contrast to quasi-PDFs, the pseudo-PDFs have  the ``canonical''  $-1 \leq x  \leq 1$ support
 for all values of $z_3^2$.  In the limit  $z_3\to 0$ they tend to PDFs,
  showing, in this  limit, a typical perturbative evolution
 with the scale $1/z_3$ being the parameter of evolution.


{As discussed in~ \cite{Radyushkin:2016hsy,Radyushkin:2017gjd},  the  fast  nonperturbative
decrease  with $z_3^2$ of   the pseudo-PDFs  $ {\cal P} (x, z_3^2) $  or the
Ioffe-time distribution  ${\cal M} (\nu, z_3^2)$,  is responsible for   delaying the approach of quasi-PDFs  $Q(y, p_3)$
to the  PDF  $f(y)$}.
 An important observation is that one can  strongly
 reduce the   $z_3^2$-dependence
by simply dividing  the Ioffe-time distribution
${\cal M} (\nu, z_3^2)$ by an appropriate  factor $D(z_3^2)$ satisfying $D(0)=1$
and having the $z_3^2$-dependence close (on average) to that of
${\cal M} (\nu, z_3^2)$.
The absence of  the \mbox{$\nu$-dependence}  in this factor and its $D(0)=1$
normalization guarantees that   the ratio ${\cal M} (\nu, z_3^2)/D(z_3^2)$
taken in the $z_3^2 \to 0$ limit  will produce the same PDF as the original function
${\cal M} (\nu, z_3^2)$ taken in the same limit.


The choice for  $D(z_3^2)$ advocated in Ref.   \cite{Radyushkin:2017cyf},
is   to take it to  be equal to  the rest-frame function $ {\cal M} (0, z_3^2)$.
An  additional advantage of this choice  is that
both ${\cal M} (\nu, z_3^2)$ and ${\cal M} (0, z_3^2)$ contain the same multiplicative
factor $Z(z_3^2)$ generated by the  renormalization of the gauge link.
In the ratio, it should cancel out.

Our  goal in the present work  is an exploratory   lattice calculation
of the $u$-$d$ proton PDF using the strategy outlined in \mbox{Ref.   \cite{Radyushkin:2017cyf}. }
 To make this article  self-contained, we reproduce in Sections 2 and 3  the main ideas
 of \mbox{Ref.   \cite{Radyushkin:2017cyf}.}   The description of the method   used for the  lattice extraction of the
 reduced Ioffe-time distribution is given in Section 4.
 The data analysis   and interpretation is discussed  in Section 5.
 The summary of the paper is given in Section 6.
